\chapter*{Conclusions and Future Work}
\markboth{Conclusions and Future Work}{}
\label{cap:conclusions}
\addcontentsline{toc}{chapter}{Conclusions and Future Work}

\section{Conclusions of the tool }

After completing the tool and reflecting on each chapter of this document, it can be concluded that a modular and functional command-line tool has been created, which aims to combine cryptography and steganography in order to communicate secretly over the public network. To justify this conclusion, we will recall the objectives defined in section \ref{sec:ob}:

\begin{itemize}
    \item \textbf{Creation of a CLI (Command-line interface)}:
    Class diagram and explanation in chapter~\ref{cap:descripcionTrabajo}.
    
    \item \textbf{Implementation of algorithms for the concealment and revelation of text in images as well as images in audio}:
    Explained in chapter~\ref{cap:descripcionTrabajo} where we can see that we have implemented algorithms based on the use of LSB (Least Significant Bit), which involves embedding text in images and transforming it to make it invisible, as well as algorithms based on the transformation of files into acoustic waves.
    
    \item \textbf{Processing of acoustic waves in computing}: 
    Explained in section~\ref{sec:audioo}.
    
    \item \textbf{Use of secret key cryptography}: 
    The fundamentals are explained in section~\ref{sec:cripto} along with its implementation in chapter~\ref{cap:descripcionTrabajo}.
\end{itemize}


\section{ Future Work }
 
Looking ahead and with the objectivity that everything can be improved, there are several aspects that would undoubtedly enhance the tool:

\begin{itemize}
    \item Implementation of error-correcting code to prevent possible man-in-the-middle attacks and noise in the transmission channel.
    
    \item Implementation of hiding algorithms for more file formats. This would provide more options for hiding.
    
    \item Possibility to add new encryptions or hiders through simple plugins for any user.
    
    \item Design of a GUI (Graphical User Interface) to be more user-friendly for the general public who may not have the knowledge to use the tool via the CLI.
    
    \item Implementation of a decoder for SSTV, inorder to avoid opening QSSTV and make the process automated.
    
    \item Implementation of a specific OCR (Optical Character Recognition) engine for a custom font and adapted to our context: black letters on a white background.
    
    \item A more powerfull analysis of \textit{ocultadorText}, generating a much larger dataset using PySpark and a cloud service like GCP (Google Cloud Console).
    
    \item Implementation of domain transformation-based techniques.
    
    \item Use of public cryptography for exchanging secret keys.

\end{itemize}


